\documentclass[output=paper]{langscibook}
\ChapterDOI{10.5281/zenodo.22880266}
\title{\textit{O paranda} – The floor}
\author{Enoch Horai Magum}
\abstract{}
\begin{document}
\maketitle\label{chap:20}
\begin{figure}[H]
    \captionsetup[subfigure]{oneside,margin={0cm,0cm}}
    \captionsetup{justification=centering}
    \subfloat[]{\includegraphics[trim=0 27 0 0,clip,width=.38\linewidth]{figures/20_1.jpg}\label{fig:20_1}}\hfil
    \subfloat[]{\includegraphics[trim=0 0 0 0,clip,width=.62\linewidth]{figures/20_2.jpg}\label{fig:20_2}}
    \caption{(a) \textit{Kobuu} -- Kobuu palm (\textit{Cyrtostachys Ioriae}),\\ (b) \textit{Paranda} -- Floor}
\end{figure}
\begin{xltabular}{\textwidth}{lXX}
    (\ref{ex:text15-1})&\textit{O paranda o peho taba vaa tea inu teara a ta Teapu.}&The floor is a thing in our house, of us, the Teop people.\\
    (\ref{ex:text15-2})&\textit{Eove to kom, hio ae goroho raara.}&It is where we step, sit, and sleep on.\\
    (\ref{ex:text15-3})&\textit{O paranda na vavaobete raara toon na naono vai to dao raara bono hogo.}&We put the floor onto the beams that we call \textit{hogo} (floor joists).\\
    (\ref{ex:text15-4})&\textit{O paranda na pakupaku raara bono kobuu.}&We make the floor from the \textit{kobuu} palm (\textit{Cyrtostachys species}).\\
    (\ref{ex:text15-5})&\textit{Tea meha maa vaan na pakupaku riori bono kumako.}&In other villages they make it from the kugumaru palm (\textit{Areca novohibernica}).\\
    (\ref{ex:text15-6})&\textit{Beara tau paku o paranda, eara toro ipinoo roho a bebeahu ae a tatabonata na inu bono kara urisi.}&When we are about to make the floor, we must measure the length and the width of the house with a tape made of the wild banana sheath.\footnote{See Chapter \ref{chap:11}:\ref{ex:text7-88}ff.}\\
    (\ref{ex:text15-7})&\textit{O kara urisi bona to kahi dee raara rasuu to kahi ipinoo raara o kobuu.}&This wild banana sheath tape we take to the jungle where we will measure the \textit{kobuu} palm.\\
    (\ref{ex:text15-8})&\textit{Beara suguna teo hum to tei nae o kobuu, eara paa sisira roho tea tabae o kobuu toro paku ara bono paranda toro vuaba vai o manoto.}&When we have arrived at the place where the \textit{kobuu} palm stands, we carefully examine it, because the \textit{kobuu} palm from which we will make the floor must be a mature one.\\
    (\ref{ex:text15-9})&\textit{Beara viana vakavara o kobuu vai tau paku ara o paranda, eara repaa mosi gunaha e.}&When we have chosen the \textit{kobuu} palm that we are about to make the floor from, we cut it down.\\
    (\ref{ex:text15-10})&\textit{O kobuu vai o kobusu saka mataa haana,}&The young and moist \textit{kobuu} palm is not good,\\
    (\ref{ex:text15-11})&\textit{eikou na antee vaha koana tea koava, evehee o vuaba vai o manoto vamataa na antee nana tea tei vananaoana.}&because it can easily rot, but the fully mature palm can last for a long time.\\
    (\ref{ex:text15-12})&\textit{Beo kobuu obete nana putaa, eara repaa ipinoo e,}&When the \textit{kobuu} palm is lying on the ground, we measure it,\\
    (\ref{ex:text15-13})&\textit{eara repaa rom suku bata koara o ipinoo.}&and then we just cut the stem crosswise according to the measure.\\
    (\ref{ex:text15-14})&\textit{Beara romrom vakavara e, eara repaa paku pete a kiu vaa tea kopaa.}&When we have finished cutting it up crosswise, we subsequently do the work of cutting it lengthwise.\\
    (\ref{ex:text15-15})&\textit{Eara he toro kopaa e teo hum buaku.}&But we must cut it lengthwise into two pieces.\\
    (\ref{ex:text15-16})&\textit{Tea kiu vaa tea kopaa, eara toro kiu nia buaku a toraaraa, ae amaa ta ginuva.}&For the work of the lengthwise cutting, we must work with two axes and some wedges.\\
    (\ref{ex:text15-17})&\textit{Tea kopaa, eara toro mosi momohu nia peha toraaraa,}&For the lengthwise cut, we must first make a cut with one axe,\\
    (\ref{ex:text15-18})&\textit{ore topo vakisi nana, beara mosi nia vaabuaku.}&so that it is still stuck (in the stem) when we make a cut with the second (axe).\\
    (\ref{ex:text15-19})&\textit{A vaabuaku re topo vakisi nana, beara rahi pute a toraaraa momohu,}&Then the second (axe) is still stuck (in the stem), when we pull out the first axe,\\
    (\ref{ex:text15-20})&\textit{beara tapa bata ni naora amaa ta ginuva teo hum to mosi vaakaa bata naona a bua toraaraa.}&while we are driving the wedges into the stem where the two axes cut it lengthwise open.\\
    (\ref{ex:text15-21})&\textit{Eara pasi paku bata koa vo naora teebona, ore antee komana beo rom kobuu takopaa.}&We will do it like this until the stem of the \textit{kobuu} palm is cut open lengthwise.\\
    (\ref{ex:text15-22})&\textit{Be kavara o hum kiu bona, eara repaa paku pete o hum vaa tea kopaa varutarutaa e.}&When this part of work is finished, we continue with the part of cutting it bit by bit lengthwise (into two halves).\\
    (\ref{ex:text15-23})&\textit{A kiu bona sa beera vaha haana,}&This work is not big again,\\
    (\ref{ex:text15-24})&\textit{na antee koana bona peha toraaraa.}&it only needs one axe.\\
    (\ref{ex:text15-25})&\textit{Beara kopaa varutarutaa bata rae, eara mee toro ipinoo batara a tatabonata bona bebeahu na kobisi kuriraara.}&While we are cutting it open lengthwise bit by bit, we must also measure the width with the length of our finger.\\
    (\ref{ex:text15-26})&\textit{Be kavara a kiu bona, eara repaa tarataraha kahie bono ogana nae}&When this work is finished, we carve its soft inside out\\
    (\ref{ex:text15-27})&\textit{ae taraha taatagi koa amaa ta kobuu, ore paa kavara kou a kiu vaa teebona.}&and just carefully carve off some pieces of \textit{kobuu} palm, and then indeed the work on it is finished.\\
    (\ref{ex:text15-28})&\textit{A kiu vaa tea paku paranda a hihivaa,}&The work of making the floor is difficult,\\
    (\ref{ex:text15-29})&\textit{eikou o kobuu o taba kikis, ei me vaaponis rakaha nana bona kiu bona.}&because the \textit{kobuu} palm is a hard thing, and this also makes this work very heavy.\\
    (\ref{ex:text15-30})&\textit{Ei, eara deedee voosu vai, are paa no vaobete e tea inu ore paa tavus paranda.}&Now, we carry it home, then we make it lie in the house, so that it becomes a floor.\footnotemark\\
\end{xltabular}\footnotetext{The floor is also made of sago-palm stems, see Chapter \ref{chap:6}:\ref{ex:text2-28}.}
\end{document}